\section{The Path Minimizer Theorem}
\label{sec:odd-path}

\begin{theorem}[Path Minimizer, Csikv\'ari and Lin \cite{csikvari2015}]
  For every odd path $P_k$ and every $n$-vertex tree $T$,
  $\Hom(P_n, P_k) \le \Hom(T, P_k)$.
\end{theorem}

\subsection{Computational verification}

To investigate the uniqueness of the path minimizer, we computationally
evaluated homomorphism counts for $H \in \{P_3, P_5, P_7, P_9\}$ and all
$n$-vertex trees with $n \le 22$. Tree enumeration uses the
Beyer--Hedetniemi canonical level-sequence algorithm~\cite{BeyerH80} implemented in C++,
producing all non-isomorphic trees. Tree counts are cross-validated
against OEIS A000055~\cite{oeis_A000055}.

\subsection{Regimes of minimality}

The eigenvalues of $P_k$ govern the degree of uniqueness with
which $P_n$ minimizes homomorphisms. We identify three regimes.

\subsubsection{Structural degeneracy}

\begin{proposition}
  \label{prop:p3}
  For any tree $T$ with bipartition $(A, B)$,
  \[
    \Hom(T, P_3) = 2^{|A|} + 2^{|B|}.
  \]
  In particular, $\Hom(P_n, P_3) = 2^{\lceil n/2 \rceil}
  + 2^{\lfloor n/2 \rfloor}$.
\end{proposition}

Since $P_n$ has the most balanced bipartition among $n$-vertex trees,
it achieves the minimum, but any tree sharing the same bipartition sizes
ties it exactly. Thus, the path minimizes the homomorphism count, but
it does not do so uniquely. This is a special case of Theorem~2.6
of~\cite{galvin2026}, since $P_3 \cong K_{1,2}$.

\subsubsection{\texorpdfstring{$P_5$}{P5}: Interior degeneracy and 2-shift invariance}

The adjacency matrix $A = A_{P_5}$ satisfies
\begin{equation}
  \label{eq:shift}
  A^3 \mathbf{1} = 3\, A \mathbf{1},
\end{equation}
where $\mathbf{1}$ is the all-ones vector. This identity implies that
the cube of the transfer matrix acts as a scalar multiple of the
matrix itself on the constant vector. Combinatorially, moving a pendant
edge two positions along the interior of a path preserves the
homomorphism count into $P_5$: the number of walks of length 3
originating from any vertex $v$ is exactly three times the number of
length-1 walks originating from $v$. Thus pendant relocations that
shift by 2 steps produce identical walk sums. This generates
exponentially growing tie classes ($\sim\!200$ ties at $n = 20$).

The eigenvalues of $P_5$ are $\pm\sqrt{3}, \pm 1, 0$, so
$\lambda_1^2 = 3$ exactly. In the language of
Lemma~\ref{lem:active_recurrence}, the active spectrum is
$\{\sqrt{3}, 0, -\sqrt{3}\}$ and the walk counts satisfy
$w_{k+2} = 3w_k$ for $k \ge 1$. The ratio
$\Hom(P_{n+2}, P_5) / \allowbreak \Hom(P_n, P_5) = 3$ for all odd $n$.

\subsubsection{Unique minimality}

For $k \ge 7$, the eigenvalues of $A_{P_k}$ are
$2\cos\bigl(\frac{j\pi}{k+1}\bigr)$ for $j = 1, \ldots, k$.
The squared spectral radius is $\lambda_1^2 = 4\cos^2\bigl(\frac{\pi}{k+1}\bigr) = 2 + 2\cos\bigl(\frac{2\pi}{k+1}\bigr)$.
By Niven's Theorem, $\cos(q\pi)$ is rational for rational $q$ if and only if
$q \in \{0, \pm 1/3, \pm 1/2, \pm 2/3, \pm 1\}$. Thus, $\lambda_1^2 \in \mathbb{Q}$ if and only if $k \in \{1, 2, 3, 5\}$.

For all $k \ge 6$, the squared dominant eigenvalue $\lambda_1^2$ is
irrational. This irrationality breaks any potential integer
shift-symmetry of the form $A^m \mathbf{1} = cA \mathbf{1}$ (which
would require $\lambda_1^m \in \mathbb{Q}$ or $\lambda_1^{m-1} \in
\mathbb{Q}$). In the computations above, this is exactly the mechanism
that separates the $P_5$ tie classes from the verified unique
minimality of $P_n$ for odd $n \ge 7$.

\begin{remark}[Irrationality]
  The verified transition to unique minimality occurs at $P_7$, the
  first odd path whose squared spectral radius ($2+\sqrt{2}$) is
  irrational. While still constructible in real radicals, that
  irrationality prevents the exact degeneracy classes seen in $P_5$.
\end{remark}

% ════════════════════════════════════════════════════════════════════════
